\chapter{מדדים שמניעים החלטות}
\label{chap:metrics}

לא כל מספר הוא מדד שכדאי לעקוב אחריו. אריק ריס מבחין בין \emph{מדדים מובילים
לפעולה} \term{(actionable metrics)} — שמראים סיבה ותוצאה ומכוונים החלטה — לבין
\emph{מדדי יוהרה} \term{(vanity metrics)} כגון מספר הורדות מצטבר, שנראים מרשימים
אך אינם משנים דבר~\cite{ries2011lean}. מדד טוב הוא בר-השוואה, מובן, ומשנה
התנהגות: אם שינוי בו אינו גורם לצוות לפעול אחרת, הוא אינו שווה מעקב.

\section{מדד-העל \texorpdfstring{\en{(North Star Metric)}}{North Star Metric}}

חברות בריאות בוחרות \emph{מדד-על} יחיד שמשקף בצורה הטובה ביותר את הערך שהמוצר
מספק ללקוח, ומיישרים סביבו את כל הצוות. בניגוד לרווח — שהוא תוצאה מאוחרת — מדד-העל
מודד את הרגע שבו הלקוח מקבל ערך אמיתי. הבחירה בו אינה טכנית אלא אסטרטגית: היא
מכריזה מהי ההצלחה שאליה החברה חותרת. טבלה~\ref{tab:northstar} מדגימה מדדי-על
אופייניים לפי סוג עסק.

\begin{table}[htbp]
  \centering
  \caption{דוגמאות למדד-העל לפי סוג העסק.}
  \label{tab:northstar}
  \begin{tabular}{l l}
    \toprule
    \textbf{סוג העסק} & \textbf{מדד-העל \en{(North Star)}} \\
    \midrule
    שוק דו-צדדי \en{(marketplace)} & עסקאות מוצלחות לשבוע \\
    \en{SaaS} צוותי               & צוותים פעילים שבועית \en{(WAU)} \\
    תוכן ומדיה                    & שעות צפייה / זמן מעורבות \\
    מסחר אלקטרוני                 & לקוחות חוזרים לרכישה חודשית \\
    \bottomrule
  \end{tabular}
\end{table}

\section{משפך \texorpdfstring{\en{AARRR}}{AARRR}: ממבקר ועד הכנסה}

מודל \en{AARRR} (``מדדי הפיראט'') מפרק את מסע הלקוח לחמישה שלבים נמדדים: רכישה
\term{(Acquisition)}, הפעלה \term{(Activation)}, שימור \term{(Retention)}, הפניה
\term{(Referral)} והכנסה \term{(Revenue)}. בכל מעבר בין שלבים יש נשירה, ואיור
\ref{fig:funnel} מציג את הירידה האופיינית. הטעות הנפוצה היא לשפוך תקציב על ראש
המשפך (רכישה) בעוד הדליפה הגדולה נמצאת בהפעלה או בשימור; עדיף לאתר את הנקודה שבה
אובדים הכי הרבה לקוחות ולתקן אותה תחילה, שכן שיפור באמצע המשפך מכפיל את עצמו על
פני כל השלבים שאחריו.

\begin{figure}[htbp]
  \centering
  \includegraphics[width=0.74\textwidth]{funnel.pdf}
  \caption{משפך ההמרה \en{AARRR} — מספר הלקוחות בכל שלב, מרכישה ועד הכנסה (גרף
  שנוצר בקוד \en{Python}). הנשירה הגדולה ביותר, ולא ראש המשפך, היא היעד לשיפור.}
  \label{fig:funnel}
\end{figure}

\section{שימור לפי קבוצות \texorpdfstring{\en{(cohort retention)}}{cohort retention}}

המדד החשוב ביותר במשפך הוא השימור, ודרך המדידה הנכונה שלו היא בקבוצות
\term{(cohorts)}: מקבצים את המשתמשים לפי שבוע ההצטרפות ועוקבים אחר האחוז שנותר
פעיל לאורך הזמן. צורת העקומה מספרת את הסיפור: עקומה שדועכת לאפס מסמנת היעדר
התאמת מוצר-שוק (פרק~\ref{chap:pmf}), בעוד עקומה ש\emph{משתטחת} ברמה חיובית מעידה
על בסיס לקוחות נאמן שעליו אפשר לבנות צמיחה (פרק~\ref{chap:growth}). השוואת קבוצות
עוקבות גם חושפת אם שינויים במוצר אכן משפרים את השימור, או רק מזיזים מספרים.

\section{מדדים מובילים מול מפגרים}

הכנסה היא מדד \emph{מפגר} \term{(lagging)}: היא מאשרת בדיעבד מה שכבר קרה. כדי
לנווט קדימה דרושים מדדים \emph{מובילים} \term{(leading)} שמנבאים את התוצאה —
למשל שיעור ההפעלה בשבוע הראשון, או מספר הפעולות המשמעותיות שמשתמש חדש מבצע. צוות
שמנהל את עצמו לפי מדדים מובילים מגלה בעיה בעוד מועד; צוות שממתין לדוח ההכנסה
החודשי מגלה אותה כשכבר מאוחר לתקן.

\section{מלכודת לוח המחוונים העמוס}

ככל שקל יותר לאסוף נתונים, כך גדל הפיתוי למלא לוח מחוונים בעשרות מדדים. התוצאה
הפוכה לכוונה: כשהכול ``חשוב'', דבר אינו חשוב, והצוות מאבד את היכולת להחליט.
משמעת מדידה טובה בוחרת מדד ראשי \emph{אחד} לכל שלב במשפך, ומתייחסת לשאר כמדדי
אבחון שנבחנים רק כאשר המדד הראשי זז לכיוון הלא נכון. חשוב לא פחות לשמור על
\emph{היגיינת נתונים}: הגדרה אחידה לכל מדד, מקור אמת יחיד, וזהירות מהטיית
ההישרדות — הסקת מסקנות מן הלקוחות שנשארו תוך התעלמות מאלה שנטשו ולכן אינם
מופיעים בנתונים.

\begin{takeaway}
בחרו \emph{מדד-על אחד} שמודד ערך אמיתי ללקוח, מדדו שימור בקבוצות (לא כממוצע
גלובלי), ונווטו לפי מדדים מובילים. מדד שאינו משנה החלטה הוא יוהרה — לא מדידה.
\end{takeaway}
